\section{三笔算力合同：从框架到部分确收}

A/B/C三单金额分别为160--190亿元、100--120亿元、130--150亿元，合计390--460亿元，均以验收后60个月按月计费。年化含税中值约85.0亿元，未税约80.2亿元。7月11日公告明确，前序合同已交付、获验收并确认收入，这是相对旧报告最重要的正向变化。

\begin{exhibitbox}[合同到盈利桥]
\noindent\begin{tabularx}{\textwidth}{L{1.5cm}R{1.8cm}L{2.1cm}X}
\toprule
\textbf{合同} & \textbf{金额} & \textbf{当前状态} & \textbf{进入模型} \\
\midrule
A & 160--190亿元 & 已生效；前序合同已验收确收 & 金额未知，不单独外推 \\
B & 100--120亿元 & 10\%预付款；前序合同已验收确收 & 金额未知，不单独外推 \\
C & 130--150亿元 & 生效，待交付、验货、验收 & 基准不直接确认 \\
合计 & 390--460亿元 & 五年合同 & 基准只确认年化未税中值25\%，净利率4\% \\
\bottomrule
\end{tabularx}
\end{exhibitbox}

\section{收入确认与净利润门槛}

基准情景把年化未税中值80.19亿元的25\%计入2026收入，按4\%净利率贡献0.8亿元；牛市按40\%确认和7.5\%净利率贡献2.4亿元。该假设已经明显低于合同稳态金额，并显式覆盖服务器折旧、融资成本和运维费用的不确定性。

\begin{riskbox}[不能把订单额直接当利润]
服务器数量、月租单价、采购成本、折旧年限、抵押融资比例、实际确认收入、毛利率和回款均未披露。三笔合同证明需求与执行开始落地，不证明390--460亿元已经成为上市公司收入。
\end{riskbox}
